% NepalDTM v2 working paper. Numbers come from validation/experiments_registry.json and the v2 validation reports;
% every value still being computed is marked \pending{...}. Compile from this directory:
%   tectonic paper_v2.tex
\documentclass[11pt]{article}

\usepackage[T1]{fontenc}
\usepackage[utf8]{inputenc}
\usepackage[margin=1in]{geometry}
\usepackage{booktabs}
\usepackage{amsmath}
\usepackage{graphicx}
\usepackage{xurl}
\usepackage{xcolor}
\usepackage[section]{placeins}  % figures and tables never float past the end of their section
\usepackage[round,authoryear]{natbib}
\usepackage[hidelinks]{hyperref}
\usepackage{fancyhdr}
\pagestyle{fancy}\fancyhf{}
\renewcommand{\headrulewidth}{0.3pt}
\fancyhead[C]{\footnotesize This is a non-peer-reviewed preprint submitted to EarthArXiv.}
\fancyfoot[C]{\thepage}
\fancypagestyle{plain}{\fancyhf{}\fancyhead[C]{\footnotesize This is a non-peer-reviewed preprint submitted to EarthArXiv.}\fancyfoot[C]{\thepage}}

\graphicspath{{figures/}}
\newcommand{\pending}[1]{\textcolor{red}{\textbf{[pending: #1]}}}
% \fig{file}{caption}{label}: shows a pending note until make_figures.py has produced the file
\newcommand{\fig}[3]{\begin{figure}[!htbp]\centering\IfFileExists{figures/#1.pdf}{\includegraphics[width=0.92\textwidth,height=0.42\textheight,keepaspectratio]{#1.pdf}}{\pending{figure \detokenize{#1}}}\caption{#2}\label{#3}\end{figure}}
\hypersetup{%
  pdftitle={NepalDTM v2: Improving on NepalDTM v1 with ICESat-2-Adjusted Stereo Strips for a National 8 m Terrain Model of Nepal},
  pdfauthor={Achyut Paudel}%
}

\title{NepalDTM v2: Improving on NepalDTM v1\\
with ICESat-2-Adjusted Stereo Strips\\
for a National 8~m Terrain Model of Nepal}
\author{Achyut Paudel\\
\small \href{mailto:achyut.pdl@gmail.com}{achyut.pdl@gmail.com}}
\date{Preprint, November 2026}

\begin{document}
\maketitle

\begin{abstract}
NepalDTM v1 corrected the NASA High Mountain Asia (HMA) 8~m stereo mosaic to bare earth with a gradient-boosting
model trained on satellite laser altimetry. Thirty-one recorded experiments showed that its error is set by the
input surface, not by the model: added predictors moved the national error by at most 2\%.
NepalDTM v2 therefore rebuilds the input. All 512 HMA stereo strips over Nepal are adjusted jointly, with a plane
per strip, against 18.5~million ICESat-2 ATL03 ground photons and 13.3~million strip-to-strip ties; and stereo
cells that v1 discarded at cliffs, where Copernicus itself is wrong, are kept wherever two adjusted strips agree.
On 19.8~million held-out ATL03 ground photons, from tracks never used in the adjustment, the v2 terrain model
scores 5.26~m RMSE (0.85~m NMAD), against 6.93~m for v1, 8.69~m for Copernicus GLO-30, 8.40~m for FABDEM and
8.97~m for GEDTM30, and it is the most accurate of the five in every land-cover, slope and elevation class.
Against an independent 2~m WorldView/GeoEye composite at Langtang the v2 surface error falls from 8.62 to 4.71~m.
We also report what did not help.
\end{abstract}

\noindent\textbf{Keywords:} NepalDTM; digital terrain model; ICESat-2 ATL03; stereo DEM strips; joint adjustment;
co-registration; Nepal; Himalaya

\section{Introduction}
Nepal has no open national terrain model finer than 30~m.
NepalDTM v1 \citep{paudel2026nepaldtm} delivered an 8~m grid by correcting the HMA 8~m stereo mosaic \citep{shean2017hma} with a
learned surface-to-ground residual, trained on ICESat-2 \citep{markus2017icesat2,neuenschwander2019atl08} and
GEDI \citep{dubayah2020gedi}.
Its national held-out error was 7.7~m RMSE on all ICESat-2 and GEDI points (the set used in
Table~\ref{tab:exps}), 6.93~m on the held-out photons used here.
This paper asks where the remaining error comes from and what can remove it with data that are openly available.

Two findings motivate v2.
First, predictors added to the residual model gave small returns: the best stack (CHMv2 canopy height
\citep{brandt2026chmv2}, survey contours and rivers, and stereo-strip agreement) reached 7.585~m against a 7.724~m
baseline on the same spatial folds (Section~\ref{sec:experiments}).
Second, the HMA strips agree with each other to 0.82~m NMAD where several overlap, yet the mosaic built from
them is offset by 1.5--2.3~m from independent references at Langtang.
The error that remains is therefore largely a co-registration error of the strips, which a geodetic adjustment
can address directly.
Co-registration of elevation data to a reference is standard in glaciology \citep{nuth2011coreg}; here all
strips are adjusted together, against ICESat-2 ATL03 ground photons and against each other.

\section{Data}
\begin{description}
\item[HMA 8~m stereo strips.] 512 along-track \citep{shean2017hmaat} and cross-track \citep{shean2017hmact} strips
over Nepal, produced with the NASA Ames Stereo Pipeline \citep{shean2016asp,shean2020hma}; 404 along-track, 109 cross-track (one strip could not be
retrieved), 17.7~GB, WGS84 ellipsoidal heights in UTM 44N and 45N.
\item[ICESat-2 ATL03 ground photons.] 547.5~million photons in 82 half-degree cells, obtained through server-side
subsetting \citep{sliderule}: ATL03 high-confidence signal photons \citep{neumann2019atl03} that the ATL08 product
\citep{neuenschwander2019atl08} labels as ground, \texttt{quality\_ph}~$=0$, and within 2~m of the median of their
10~m along-track bin. Heights are converted to EGM2008 \citep{pavlis2012egm2008}. Whole reference ground tracks
($\approx$20\%) are held out by a fixed hash of the track number and never used for fitting.
\item[Independent 2~m references.] The gated median of 12 WorldView/GeoEye stereo pairs at Langtang Lirung,
2015--2017 \citep{shean2026langtang}, and a 2~m composite of 38 cross-track pairs at Chamoli, India, September 2015
\citep{bhushan2021chamoli}. Both are surface models, converted from ellipsoidal heights to EGM2008, with pixels
of low pair count or NMAD $\geq 3$~m removed.
\item[Ancillary.] Copernicus GLO-30 \citep{esa2021copernicus}, FABDEM \citep{hawker2022fabdem} (void-fill
comparison only), ESA WorldCover \citep{zanaga2022worldcover},
CHMv1 \citep{tolan2024canopy} and CHMv2 canopy height \citep{brandt2026chmv2}.
\item[Ground control.] Differential GPS on and beside Khumbu Glacier, 2015--2016 \citep{quincey2016khumbu}.
\end{description}

\section{Method}
\subsection{Joint strip adjustment}
For strip $k$ with centre $(x_k, y_k)$ we estimate a plane correction
\begin{equation}
c_k(x,y) = a_k + b_k\,\frac{x-x_k}{10\,\text{km}} + d_k\,\frac{y-y_k}{10\,\text{km}} .
\end{equation}
Observations are gathered per 24~km tile core of the national grid, on stable open ground (WorldCover not tree,
snow or water; slope $<35^\circ$), after removing blunders more than 40~m from the per-pixel median of all strips:
\begin{align}
z_k(x,y) + c_k(x,y) - h_{\text{photon}} &= 0, \\
\bigl(z_k + c_k\bigr) - \bigl(z_j + c_j\bigr) &= 0 \quad \text{where strips $k$ and $j$ overlap},
\end{align}
plus a weak prior pulling all parameters to zero so that strips without photons stay defined.
Tie points more than 25~m from the median for a strip or strip pair are discarded.
The sparse system ($3\times512$ unknowns) is solved with LSQR \citep{paige1982lsqr} in SciPy \citep{virtanen2020scipy}
by iteratively reweighted least squares with Huber weights \citep{huber1964robust}.
Planes are centred on each strip, so one correction holds in every tile the strip crosses.
The adjusted surface is the per-pixel median of the corrected strips, with strip count and NMAD kept as layers.

\subsection{Terrain model v2}
The adjusted median replaces the mosaic as the stereo surface.
A blunder filter rejects adjusted cells more than 40~m from the local median difference to Copernicus, except
where at least two strips agree to better than 5~m NMAD (Section~\ref{sec:voids}).
Remaining gaps are filled from Copernicus with the local offset spread by normalized Gaussian convolution
\citep{knutsson1993normconv}, as in v1.
The residual model of v1, a histogram gradient-boosting regressor \citep{pedregosa2011sklearn,ke2017lightgbm},
is retrained with three added predictors: strip count, strip NMAD and CHMv2 canopy
height.

\subsection{Metrics}
The headline check is the DTM against held-out ICESat-2 ATL03 ground photons: whole ground tracks never used in
the adjustment or the model, every model scored on the same photons.
The v1 protocol is also reported for continuity: GroupKFold over 48~km spatial blocks, scored on held-out ATL08
segments, whose 100~m fits understate gains at 8~m.
The Langtang and Chamoli composites and Khumbu dGPS test the surface model (DEM).
RMSE is the root-mean-square error; NMAD is 1.4826 times the median absolute deviation from the median.

\section{Results}
\subsection{Langtang test area}
Table~\ref{tab:langtang} compares the v1 surface, the unadjusted strip median and the adjusted strip median
(26 strips; 1.26~M photon ties, 0.70~M overlap ties).
\begin{table}[!htbp]
\centering
\caption{Langtang Lirung test area. RMSE (NMAD) in metres. Cells where both surfaces have stereo only (voids: Section~\ref{sec:voids}). Left: against the gated 2~m composite averaged to 8~m
(1,186,973 cells). Right: against held-out ICESat-2 ground photons on open ground (55,975 photons).}
\label{tab:langtang}
\begin{tabular}{lccc|ccc}
\toprule
 & \multicolumn{3}{c|}{2~m reference} & \multicolumn{3}{c}{held-out photons} \\
Slope & v1 & strip median & adjusted & v1 & strip median & adjusted \\
\midrule
all       & 8.41 (1.28) & 6.23 (1.20) & \textbf{5.88} (0.92) & 4.03 (1.40) & 3.61 (1.84) & \textbf{2.61} (1.15) \\
$<20^\circ$  & 3.03 (0.92) & 2.57 (1.01) & \textbf{1.97} (0.50) & 2.72 (1.01) & 2.51 (1.55) & \textbf{1.30} (0.78) \\
$20$--$40^\circ$ & 5.63 (1.09) & 3.84 (1.01) & \textbf{3.38} (0.83) & 4.52 (2.53) & 4.26 (2.39) & \textbf{3.18} (2.38) \\
$\geq40^\circ$ & 12.46 (2.02) & 9.38 (1.82) & \textbf{9.05} (1.66) & 10.97 (4.92) & 8.85 (4.96) & \textbf{7.76} (4.82) \\
\bottomrule
\end{tabular}
\end{table}
The adjustment also removes the systematic offset: the median difference against the 2~m reference moves from
$+1.60$ to $-0.40$~m, and against the photons from $+2.35$ to $+0.20$~m.

\subsection{National solve}
The national system has 18,506,558 photon ties and 13,320,995 overlap ties for 512 strips, of which 443 have
photon ties and 69 are held only through overlaps.
The Huber solve converges after four iterations to 1.22~m NMAD for photon ties and 1.17~m for overlap ties.
Strip offsets have a median of $-1.37$~m and a 5--95\% range of $-6.44$ to $+3.23$~m (Figure~\ref{fig:strips}).
\fig{fig01_strip_corrections}{National solve: vertical offset (left) and tilt magnitude (right) of the 512 strip
planes.}{fig:strips}

\subsection{Where the remaining surface error sits: voids}
\label{sec:voids}
Table~\ref{tab:langtang} scores only cells where both surfaces have stereo.
Over every 8~m cell of the Langtang reference, a first v2 product scored no better than v1 (8.67 against 8.62~m),
although the adjusted strips passed into it unchanged.
Splitting by source explains it (Figure~\ref{fig:source}): on the 94\% of cells with accepted stereo the error
fell from 3.06 to 2.66~m and the $+1.6$~m bias disappeared, while the 6\% filled from Copernicus scored 32~m and
dominated the total.
The voids are mostly cliffs and shadowed faces, exactly where Copernicus itself is wrong (Figure~\ref{fig:maps}).
\fig{fig03_langtang_error_maps}{DEM minus the 2~m composite over Langtang Lirung, v1 (left) and v2
(right). Red: DEM above the reference.}{fig:maps}
\fig{fig04_error_by_source}{Langtang DEM RMSE by source of the 8~m cell, v1 and v2.}{fig:source}
The v1 blunder filter compared every stereo cell with Copernicus and so also rejected good stereo there.
We compared five fills on the Langtang void cells (EXP030, Table~\ref{tab:voids}).
Keeping adjusted cells where at least two strips agree (NMAD $<5$~m) is the only arm that helps; FABDEM, a wider
offset kernel and photon-anchored offsets change nothing.
\begin{table}[!htbp]
\centering
\caption{Void filling at Langtang (EXP030). RMSE (NMAD) in metres against the 2~m composite on void cells
(80,549), and on the subset where re-admitted strips exist (30,747).}
\label{tab:voids}
\begin{tabular}{lcc}
\toprule
Fill & all void cells & re-admitted cells \\
\midrule
Copernicus + offset (v1) & 32.72 (25.90) & 47.83 \\
Copernicus + wide offset & 36.83 (27.41) & 54.03 \\
FABDEM + offset & 32.82 (26.21) & 47.92 \\
photon-anchored offset & 32.69 (25.88) & 47.79 \\
\textbf{re-admit agreed strips} & \textbf{14.25} (4.12) & \textbf{3.75} (1.17) \\
\bottomrule
\end{tabular}
\end{table}
The national stack was rebuilt with this rule; it is the v2 reported below.

\subsection{National held-out accuracy}
Table~\ref{tab:photons} and Figure~\ref{fig:photons} give the headline check: 19.8~million held-out ATL03 ground
photons (a 25\% sample of the held-out tracks), every model scored on the same photons.
v2 lowers the error in every class, most where v1 relied on Copernicus fill: slopes of $40^\circ$ and
more from 29.62 to 21.39~m RMSE, bare ground and snow from 9.05 to 6.52~m, the high Himal from 8.55 to 6.10~m.
On the v1 protocol (Table~\ref{tab:national}) the gain is small, as expected from a target smoothed over 100~m.
\begin{table}[!htbp]
\centering
\caption{v1 protocol: spatial-block held-out RMSE in metres against ICESat-2 ATL08 (6.24~M points).}
\label{tab:national}
\begin{tabular}{lrr}
\toprule
Subset & v1 & v2 \\
\midrule
all & 5.49 & \textbf{5.48} \\
forest & 6.82 & \textbf{6.63} \\
slope $<20^\circ$ & 3.26 & \textbf{3.23} \\
slope $\geq40^\circ$ & 10.80 & \textbf{10.57} \\
\bottomrule
\end{tabular}
\end{table}
Over every 8~m cell of the Langtang composite, voids included, the v2 DEM error is 4.71~m against 8.62~m for v1
(2.65 against 4.69~m under $20^\circ$; Figure~\ref{fig:maps}).
On Khumbu dGPS the v2 DEM scores 2.47~m at the glacier margins, 2.51~m on the glacier surface and 1.19~m on the
wind-tower grid (v1: 2.49, 2.96 and 1.02~m).
\fig{fig08_heldout_photons}{DTM error of Copernicus GLO-30, FABDEM, GEDTM30 and NepalDTM v1 and v2 against
19.8~M held-out ICESat-2 ATL03 ground photons (tracks never used in the adjustment): RMSE (top) and NMAD
(bottom).}{fig:photons}

\subsection{Comparison with public global DEMs}
The widely used open alternatives for Nepal are Copernicus GLO-30 \citep{esa2021copernicus}, a 30~m surface
model; FABDEM \citep{hawker2022fabdem}, Copernicus with forests and buildings removed by machine learning; and
GEDTM30 \citep{ho2025gedtm30}, a 30~m terrain model trained partly on ICESat-2, which makes its photon scores
optimistic.
All three are on EGM2008, as is NepalDTM, and are resampled bilinearly to the same 8~m grid.
Each is scored on exactly the cells and points used for NepalDTM: held-out ATL03 photon tracks (DTM), the Langtang
2~m composite (surface models only) and Khumbu dGPS (Figures~\ref{fig:photons} and~\ref{fig:public}).
\begin{table}[!htbp]
\centering
\caption{DTM error against 19.8~M held-out ICESat-2 ATL03 ground photons, all models on the same photons.
RMSE (NMAD) in metres. GEDTM30 was trained partly on ICESat-2.}
\label{tab:photons}
\small
\begin{tabular}{lrrrrr}
\toprule
Group & Copernicus & FABDEM & GEDTM30 & v1 & \textbf{v2} \\
\midrule
all & 8.69 (1.14) & 8.40 (1.11) & 8.97 (1.83) & 6.93 (0.93) & \textbf{5.26} (0.85) \\
forest & 10.93 (4.19) & 8.56 (2.86) & 8.24 (2.90) & 6.95 (2.37) & \textbf{6.71} (2.21) \\
non-forest & 8.39 (0.97) & 8.38 (0.99) & 9.05 (1.73) & 6.92 (0.83) & \textbf{5.07} (0.76) \\
bare / snow & 10.92 (1.93) & 10.92 (1.98) & 11.85 (3.36) & 9.05 (1.62) & \textbf{6.52} (1.43) \\
slope $<20^\circ$ & 4.80 (0.79) & 4.54 (0.80) & 4.93 (1.36) & 3.17 (0.65) & \textbf{2.80} (0.60) \\
slope $20$--$40^\circ$ & 11.86 (3.85) & 11.38 (3.68) & 12.26 (5.41) & 9.40 (3.01) & \textbf{7.04} (2.83) \\
slope $\geq40^\circ$ & 34.02 (11.47) & 33.49 (10.81) & 35.17 (13.42) & 29.62 (7.49) & \textbf{21.39} (6.58) \\
Terai $<300$~m & 2.68 (0.47) & 1.81 (0.46) & 1.54 (0.58) & 1.37 (0.38) & \textbf{1.34} (0.38) \\
Siwalik & 9.13 (3.16) & 7.54 (2.77) & 7.05 (2.75) & 5.93 (1.88) & \textbf{5.73} (1.61) \\
hills 1--2.5~km & 10.87 (4.35) & 9.71 (4.41) & 9.70 (4.84) & 8.29 (3.33) & \textbf{7.90} (3.10) \\
mountains 2.5--4~km & 18.41 (6.88) & 16.49 (5.77) & 16.52 (6.56) & 13.27 (4.11) & \textbf{12.05} (3.58) \\
high Himal $>4$~km & 10.31 (1.67) & 10.32 (1.72) & 11.23 (3.11) & 8.55 (1.45) & \textbf{6.10} (1.28) \\
\bottomrule
\end{tabular}
\end{table}
Against the best public DEM in each class, v2 lowers the RMSE by 37\% overall (5.26 against 8.40~m, FABDEM),
19\% under forest (6.71 against 8.24~m, GEDTM30), 38\% on slopes under $20^\circ$ and 36\% on slopes of
$40^\circ$ and more.
Over the Langtang composite, Copernicus scores 13.96~m, FABDEM 14.34~m and GEDTM30 17.37~m against 4.71~m for
the v2 DEM; on Khumbu dGPS Copernicus scores 3.87, 4.30 and 1.52~m (margins, glacier surface, wind tower) against
2.47, 2.51 and 1.19~m.
The forest removal of FABDEM and GEDTM30 lowers the error under trees relative to Copernicus, but both score
worse than Copernicus on the Langtang surface.
On the national laser set inside Nepal (4.47~M ATL08 points, spatial blocks; v2 validation report), v2 scores
6.16~m against 9.62~m for Copernicus, 9.09~m for GEDTM30 and 8.40~m for FABDEM.
\fig{fig11_public_dems}{Public global DEMs against NepalDTM on independent references: the Langtang 2~m composite
by slope (left) and Khumbu dGPS (right).}{fig:public}

\subsection{Bare ground in the DTM}
At the Khumbu dGPS sites \citep{quincey2016khumbu} the v2 DTM lies 6--8~m below the DEM on bare moraine and
glacier, and keeping the DEM on WorldCover bare, snow/ice and moss/lichen cells would restore the DEM accuracy
there (wind tower 7.63 to 1.19~m RMSE; EXP031).
On the 8.9~million held-out photons on bare ground and snow, however, that rule raises the RMSE from 6.52 to
7.49~m (all photons: 5.26 to 5.81~m), so the residual model corrects a real surface bias over snow and rock in
general and Khumbu is a local exception.
The rule is rejected; at Khumbu users should prefer the DEM.

\section{Experiments that did not change the national picture}
\label{sec:experiments}
All experiments are recorded in \texttt{validation/experiments\_registry.json} of the code repository, with the same spatial folds unless noted.
\begin{table}[!htbp]
\centering
\caption{Selected experiments. RMSE in metres on the national laser set (13.27~M points) unless noted.}
\label{tab:exps}
\begin{tabular}{llrr}
\toprule
ID & Change & baseline & result \\
\midrule
EXP002 & CHMv2 canopy height as predictor & 7.722 & 7.683 \\
EXP003 & subtract CHMv2 from the target & 7.721 & 7.981 \\
EXP018 & add GEDI to ICESat-2 training (ICESat-2 scored) & 5.402 & 5.616 \\
EXP019 & force ground downhill along surveyed rivers (near rivers) & 8.179 & 34.168 \\
EXP007--010 & physics proxies: land use, snow, phenology, soil carbon & 7.722 & $\pm0.00$ \\
EXP021 & separate models for forest, cropland, other & 7.721 & 7.700 \\
EXP022 & CHMv2 + contour + river predictors & 7.721 & 7.641 \\
EXP024 & EXP022 + stereo-strip count, NMAD, span & 7.724 & 7.585 \\
EXP028 & smooth per-strip correction (4~km knots), Langtang 2~m & 5.88 & 6.00 \\
EXP029 & train on 8~m ATL03 photon cells (held-out photon cells) & 4.14 & 4.11 \\
EXP030 & re-admit agreed strips in voids (Langtang void cells) & 32.72 & \textbf{14.25} \\
EXP032 & Sentinel-1 72-pair InSAR height correction of Copernicus, Langtang 2~m & 12.70 & 12.69 \\
-- & GEM-15 DSM (bias removed) as fill, held-out photons on fill cells & 3.67 & 10.64 \\
\bottomrule
\end{tabular}
\end{table}
Proxies built from elevation, slope, aspect and land cover add no information the model does not already have
(Figure~\ref{fig:exps}).
\fig{fig10_experiments}{Change in RMSE of each experiment against its own baseline. Only the two changes to the
input surface (EXP027, EXP030) move the error by more than a few percent.}{fig:exps}
Subtracting canopy height from a 10~m product at 8~m introduces registration noise.
Hydrological enforcement fails because surveyed rivers carry no heights, so the downhill direction must come from the
surface being corrected.

Two further sources were tested as replacements for the Copernicus fill.
A multi-baseline Sentinel-1 stack (72 HyP3 interferograms~\citep{hyp3} over three tracks, time gaps of at most
36~days) was inverted per pixel for the height error of Copernicus, which HyP3 uses as the reference topography.
The solved correction has a median magnitude of 0.11~m and leaves the Copernicus error unchanged (12.70 to
12.69~m RMSE against the Langtang 2~m composite; EXP032): with perpendicular baselines of 40--60~m (median), a 10~m height
error produces about 0.2~rad of phase, below the atmospheric and decorrelation noise.
GEM-15~\citep{gem15}, a 15~m DSM and DTM from sub-metre commercial stereo, is better than Copernicus at Langtang
(6.40 against 11.59~m RMSE against the 2~m composite) but worse than v2 (4.05~m).
On 9.0~M held-out photons on the cells where v2 uses Copernicus fill, GEM-15 is about 14~m too high (median), as
expected from the deliberate offset its producers apply over India and taper across the border; removing a per-tile
bias leaves 10.64~m RMSE against 3.67~m for the v2 DTM there.
Only on slopes of $40^\circ$ and more does the debiased GEM-15 surface (12.00~m) edge out the v2 DTM (12.75~m).

\subsection{Transfer outside Nepal}
Run unchanged over Chamoli (Uttarakhand), outside the training area, the v1 pipeline gives a surface error of
9.22~m RMSE against the 2~m composite, versus 14.47~m for Copernicus (3.68 versus 6.55~m under $20^\circ$).
That composite was co-registered to an HMA mosaic, so part of the surface agreement is shared lineage.

\section{Discussion and limitations}
The adjustment uses a plane per strip. A smoother correction surface fits the photons slightly better (1.41 to
1.35~m NMAD) but scores worse against the 2~m reference (EXP028), so it overfits the sparse ties.
Training on 8~m photon cells instead of ATL08 segments gains 0.7\% (EXP029): the residual model is limited by
what its predictors explain, not by the smoothness of its target.
Photon ties exist only on open, gentle ground; forest and slopes above $35^\circ$ inherit the strip correction
estimated elsewhere on the strip.
Steep forest remains the hardest class: none of the open data sees the ground under closed canopy.
Only 39\% of Nepal's area has accepted HMA stereo; the rest still relies on Copernicus fill, which carries most
of the remaining surface error (Section~\ref{sec:voids}).
The 2~m references are surface models and cover two small areas; they test the surface, not the bare earth.
A single RMSE over all cells hides the split between stereo and fill; we report both.

\section{Data and code availability}
Code: \url{https://github.com/achyut/NepalDTM}, branch \texttt{v2} (\texttt{code/run\_v2.sh} runs the full build;
\texttt{code/strip\_adjust.py}, \texttt{code/nw\_strip\_adjust.py}, \texttt{code/fetch\_atl03.py}; experiment records in
\texttt{validation/experiments\_registry.json}).
Data: the v2 DEM, DTM and auxiliary layers as 324 Cloud-Optimized GeoTIFF tiles with VRT mosaics, and the DTM as
Terrarium-encoded PMTiles for web maps, archived on Zenodo,
\url{https://doi.org/10.5281/zenodo.24327134}~\citep{paudel2026nepaldtmv2data} (v1: \url{https://doi.org/10.5281/zenodo.24327066}~\citep{paudel2026nepaldtmv1data}),
and mirrored for streaming at \url{https://nepaldtm-data.achyutpdl.com/v2/}; interactive viewer at
\url{https://nepaldtm.achyutpdl.com}.
The Langtang and Chamoli composites are CC~BY-NC~4.0; derived products that include them inherit that licence.

\FloatBarrier
\bibliographystyle{plainnat}
\bibliography{references}

\end{document}
